% 第8.5节；首次阅读翻译，按用户要求不进行额外复核。
\section{Seidel 关于四次 K3 曲面的结果}\label{sec:8.5}
\setcounter{equation}{26}

\SourceParagraph{ch08-8.5-001}{661}
在最后一节，简要讨论 Paul Seidel\BookCite{ref420}的工作：他在四次 K3 曲面这一极不平凡的情形中证明了同调镜像对称。准确地说，设 $X$ 为 $\mathbb P^3$ 中的非奇异四次曲面，配备适当的辛形式，例如 Fubini--Study 度量的 Kähler 形式在 $X$ 上的限制。另一方面，考虑镜像族；这里关注的是它的复结构。取 $\mathbb P^3\times\operatorname{Spec}\mathbb C[[q]]$ 中由
\begin{equation}\label{eq:8.27}
q(x_0^4+x_1^4+x_2^4+x_3^4)+x_0x_1x_2x_3=0
\end{equation}
定义的族，再对下列群的作用取商：
\[
\Gamma_{16}:=
\left\{(a_0,a_1,a_2,a_3)\ \middle|\
 a_i\in\mathbb C,\ a_i^4=1,\ \prod_{i=0}^3a_i=1\right\}
\big/\left\{(a,a,a,a)\mid a^4=1\right\}.
\]
它在 $\mathbb P^3$ 上以对角方式作用，即 $(a_0,a_1,a_2,a_3)$ 将
\[
(x_0,x_1,x_2,x_3)\longmapsto
(a_0x_0,a_1x_1,a_2x_2,a_3x_3).
\]
将上述族对 $\Gamma_{16}$ 的商记为 $\check X\to\operatorname{Spec}\mathbb C[[q]]$。其一般纤维的有理双点经过保典消解，得到域 $\mathbb C((q))$ 上的 K3 曲面 $\widehat X_q$，因而可以考虑范畴 $D^b(\widehat X_q)$，即 $\widehat X_q$ 上凝聚层的有界导出范畴。可将此范畴与 $\Lambda_{\mathrm{nov}}$ 作张量积；后者包含 $\mathbb C((q))$ 的代数闭包。也就是说，将所有态射空间都与 $\Lambda_{\mathrm{nov}}$ 作张量积。Seidel 的结果大体上可表述如下。（先声明一下：本书给出的定义，与 Seidel 论文\BookCite{ref420}所需的定义，在技术上有一些差异。精确陈述请见\BookCite{ref420}。）

\SourceParagraph{ch08-8.5-002}{662}
\BookStatementNumber{33}
\begin{theorem}\label{thm:8.33}
存在 $\psi\in\operatorname{Aut}(\mathbb C[[q]])$，以及三角 $A_\infty$ 范畴之间的等价
\[
D^\pi\operatorname{Fuk}(X)
\cong\psi^*D^b(\widehat X_q)\otimes\Lambda_{\mathrm{nov}}.
\]
\end{theorem}

\SourceParagraph{ch08-8.5-003}{662}
应将 $\psi$ 看作镜像映射，它在两侧的 $q$ 之间变换坐标。因为 $q$ 不是族 $\check X\to\operatorname{Spec}\mathbb C[[q]]$ 的典范坐标，所以预计这一变换并不平凡。不过，目前还不知道这个坐标变换 $\psi$ 是否确实对应于典范坐标。

\SourceParagraph{ch08-8.5-004}{662}
Seidel 的证明与这里给出的椭圆曲线情形的证明很不一样，而且相当复杂。这里只能满足于一个极简略的证明思路。

\SourceParagraph{ch08-8.5-005}{662}
基本想法是在两侧分别找到一小组对象，它们分裂生成相应的范畴。若通过逐次取映射锥、分裂出直和项以及取同构，能够从一组对象得到范畴 $\mathcal A$ 中的任意对象，就称这组对象分裂生成 $\mathcal A$。这些对象所生成的范畴，可以用显式方法研究并加以比较。

\SourceParagraph{ch08-8.5-006}{662}
在凝聚层一侧，Beilinson 定理，即\BookRef{thm:4.72}{定理4.72}，说明 $D^b(\mathbb P^3)$ 由一小组局部自由层生成，即 $\Omega^i(i)[i]$，$0\leq i\leq3$。事实上，这些层在四次 K3 曲面上的限制，分裂生成该曲面的有界导出范畴。在我们的情形中，$\check X$ 被描述为一个 K3 曲面族 $\mathcal X$ 对 $\Gamma_{16}$ 的商。因而，利用导出 McKay 对应，可以用商空间保典消解的导出范畴所对应的、覆盖空间上的 $\Gamma_{16}$ 等变导出范畴来研究问题。覆盖空间上的每个分裂生成元都有 16 个等变版本，从而得到 $D^b(\widehat X_q)$ 的一组 64 个分裂生成元。

\SourceParagraph{ch08-8.5-007}{662}
在另一侧，所考虑的 Fermat 族\BookEqRef{eq:8.27}{8.27}，作为 $\operatorname{Spec}\mathbb C[q]$ 上的族，除了 $q=0$ 处的退化纤维外，还有四个奇异纤维，每个纤维都有 16 个奇异点。每个这样的奇异点，都会在此族中某个选定的一般纤维内诱导一个作为消灭闭链的 Lagrange 球面。这样便在该纤维中得到 64 个 Lagrange 球面；适当地选择它们，就得到一组对象，用来与导出范畴中的对象比较。可以证明，这 64 个 Lagrange 子流形分裂生成 $D^\pi\operatorname{Fuk}(X)$。最精细的一步，是计算这些对象所生成的子范畴；这需要大量计算，特别需要计算若干乘法映射。最终，由此可以比较这两个范畴。

\SourceParagraph{ch08-8.5-008}{663}
事实上，可以证明，两侧的范畴都是同一个范畴的非平凡变形。在两侧，$q$ 都是变形参数。通过 Hochschild 上同调的计算，关于 Hochschild 上同调的一些讨论见\BookSectionRef{sec:2.2.3}{2.2.3}，能够证明，在所考虑的对称性约束下，只有一个单参数变形族。因此，这两个族必定相同，但无法给出显式的镜像映射。
